% AI-produced by GPT-6 Astra, OpenAI, 20 November 2026.
% Unnumbered gate diagnostic; not an A5 counterexample or gate closure.
\subsection*{A fixed Hardt partition does not detect the exact bundle locus}

\textbf{Claim.}\quad
Over every real closed field there is a continuous proper semialgebraic
map $f:X\to R^3$ with every fibre semialgebraically homeomorphic to the
same closed PL two-manifold $F$, and a Hardt triviality partition
\[
 R^3=S\sqcup(R^3\setminus S),\qquad S=R(1,0,0),
\tag{HT.1}
\]
such that its exact definable bundle locus is $R^3\setminus\{0\}$.
In particular this locus is not a union of the connected pieces of the
displayed Hardt partition. The example has a definable bundle locus;
it does not contradict A5 or exclude a richer finite transverse-germ
recognition theorem.

\begin{proof}
\textbf{The circle, its folded map, and the parameter homotopy.}\quad
Let $C=\partial[0,1]^2$. Define the piecewise-affine parametrization
$c:[0,1]\to C$ by
\[
 c(s)=
 \begin{cases}
 (1,4s),&0\le s\le1/4,\\
 (2-4s,1),&1/4\le s\le1/2,\\
 (0,3-4s),&1/2\le s\le3/4,\\
 (4s-3,0),&3/4\le s\le1.
 \end{cases}
\tag{HT.2}
\]
The formulas agree at their common endpoints. The only repeated point
is $c(0)=c(1)=(1,0)$; otherwise $c$ is injective. Put
\[
 b(s)=
 \begin{cases}
 3s,&0\le s\le1/3,\\
 2-3s,&1/3\le s\le2/3,\\
 3s-2,&2/3\le s\le1.
 \end{cases}
 \qquad
 a_\lambda(s)=(1-\lambda)s+\lambda b(s),\quad0\le\lambda\le1.
\tag{HT.3}
\]
This is a jointly continuous semialgebraic function into $[0,1]$,
with $a_\lambda(0)=0$ and $a_\lambda(1)=1$. Hence the formula
\[
 q_\lambda(c(s))=c(a_\lambda(s))
\tag{HT.4}
\]
is well-defined, continuous jointly in $(\lambda,c(s))$, and
semialgebraic. Here is the endpoint justification without a hidden
quotient assumption. Off $(1,0)$ one can use the continuous inverse
of $c$ on the corresponding open circle arc. When $c(s)$ approaches
$(1,0)$, its representative approaches either $0$ or $1$; continuity
of $a_\lambda$ uniformly for $\lambda$ in its compact interval makes
the corresponding image approach $c(0)=c(1)$. This proves the joint
continuity at the identified endpoint as well. The definable
intermediate value theorem applied to $a_\lambda$ shows that every
$q_\lambda$ is onto. Moreover $q_0=\operatorname{id}_C$.

Let $F=C\times C\subset R^4$ and
$Q_\lambda=q_\lambda\times\operatorname{id}_C$. Thus $F$ is a
fixed closed PL torus, each $Q_\lambda$ is a continuous semialgebraic
surjection, and $Q_0=\operatorname{id}_F$.

\textbf{An exact obstruction to approximating $Q_1$ by homeomorphisms.}\quad
Choose
\[
 p=(c(1/8),c(1/8))\in F.
\]
Since $c(1/8)$ is not the identified endpoint, its three inverse
representatives under $q_1$ are exactly the solutions of $b(s)=1/8$:
\[
 s_1=1/24,\qquad s_2=5/8,\qquad s_3=17/24.
\tag{HT.5}
\]
They lie strictly inside the three affine pieces in (HT.3).
It follows that
\[
 Q_1^{-1}(p)=\{x_1,x_2,x_3\},
 \qquad x_i=(c(s_i),c(1/8)),
\tag{HT.6}
\]
with three distinct points.

We prove directly that there is $\varepsilon_0>0$ for which no
definable homeomorphism $H:F\to F$ satisfies
$\sup_x\|H(x)-Q_1(x)\|_2<\varepsilon_0$.
Take pairwise disjoint relative open neighborhoods $O_i$ of $x_i$,
with pairwise disjoint closures, small enough that their union does
not cover $F$. The complement
$K=F\setminus\bigcup_i O_i$ is definably compact and contains no
inverse image of $p$. Thus
\[
 d_0=\min_{x\in K}\|Q_1(x)-p\|_2>0.
\tag{HT.7}
\]
Both coordinates of $p$ lie in the interiors of edges of the square
circle. There is consequently $r_0>0$ such that
$F\cap B_{r_0}(p)$ is the intersection of that ball with one affine
two-plane, hence is convex. Choose
$0<\delta<\min\{r_0,d_0/3\}$ and put $\varepsilon_0=\delta$.

If $H$ were that close, $H(x_1)$ and $H(x_2)$ would lie in the
convex set $F\cap B_\delta(p)$. Join them by its straight segment
$\sigma:[0,1]\to F\cap B_\delta(p)$, and let
$\gamma=H^{-1}\sigma$. For every $s$,
\[
 \|Q_1(\gamma(s))-p\|_2
 \le\|Q_1(\gamma(s))-H(\gamma(s))\|_2
       +\|\sigma(s)-p\|_2
 <2\delta<d_0.
\]
Hence $\gamma$ takes values in $\bigcup_i O_i$ by (HT.7), but
joins $x_1\in O_1$ to $x_2\in O_2$. This contradicts definable
connectedness of the interval: the inverse images of the finitely
many disjoint relative open sets form a nontrivial clopen partition.
Over the real numbers the same proof applies to arbitrary topological
homeomorphisms using ordinary connectedness. This establishes the
obstruction completely, without invoking a general theorem about
limits of homeomorphisms.

\textbf{The proper family over a three-dimensional base.}\quad
For $v\in R^3\setminus\{0\}$ write
$t=\|v\|_\infty>0$ and $u=v/t\in\partial[-1,1]^3$.
Define the semialgebraic subset of $R^3\times R^4\times R^4$
\[
\begin{split}
 X={}&\{(v,t x,Q_{|u_2|}(x)):
                v\ne0,\ t=\|v\|_\infty,\ u=v/t,\ x\in F\}\\
 &\quad\cup\{(0,0,y):y\in F\},
 \qquad f(v,z,y)=v.
\end{split}
\tag{HT.8}
\]
The set $X$ is closed. Away from $v=0$ it is the graph of the
displayed continuous function after the coordinate recovery $x=z/t$.
At a point in its closure with $v=0$, boundedness of $x\in F$
forces $z=0$, and closedness of $F$ forces $y\in F$; such points
are included in (HT.8). This also proves closedness over an arbitrary
real closed field, either directly from neighborhoods or by definable
curve selection. If the base coordinates are bounded, then $t$, $t x$
and $y\in F$ are bounded. Therefore the inverse image of every closed
bounded definable set in the base is closed and bounded. This proves
definable properness. Over $R=\mathbb R$ it proves ordinary properness.

Each noncentral fibre is homeomorphic to $F$ through its variable $x$,
and the central fibre is $F$ through $y$. Off the origin, the explicit
homeomorphism
\[
 (x,v)\longmapsto(v,\|v\|_\infty x,Q_{|v_2|/\|v\|_\infty}(x))
\tag{HT.9}
\]
trivializes $f$, with inverse recovering $x$ by division by
$\|v\|_\infty$. On the whole axis $S$, including its origin,
$|u_2|=0$ and $Q_0=\operatorname{id}_F$. The map
\[
 F\times R\longrightarrow f^{-1}(S),\qquad
 (x,a)\longmapsto(ae_1,|a|x,x)
\tag{HT.10}
\]
is an exact product homeomorphism, whose inverse uses the last
coordinate and the first base coordinate. Equations (HT.9)--(HT.10)
prove Hardt triviality over both pieces in (HT.1). Both pieces are
connected definable manifolds.

\textbf{Failure only at the origin.}\quad
Suppose a definable product chart existed near $0$. Using its central
fibre homeomorphism to change the fibre coordinate, write it as
$H:F\times V\to f^{-1}(V)$ with
$H(z,0)=(0,0,z)$ and $fH(z,v)=v$. For small positive $t$ restrict
to $v=te_2$. Equation (HT.8) gives homeomorphisms $h_t:F\to F$
with
\[
 H(z,te_2)=(te_2,t h_t(z),Q_1(h_t(z))).
\tag{HT.11}
\]
Continuity on the definably compact set $F$ times a small closed base
cube implies uniformly in $z$ that
$Q_1(h_t(z))\to z$ as $t\downarrow0$. Since each $h_t$ is onto,
\[
 \sup_{x\in F}\|Q_1(x)-h_t^{-1}(x)\|_2
 =\sup_{z\in F}\|Q_1(h_t(z))-z\|_2\longrightarrow0.
\tag{HT.12}
\]
This contradicts the positive lower bound $\varepsilon_0$ just
proved. Hence $0$ is not a bundle value. Equation (HT.9) proves
that every other value is. This establishes the claim.
\end{proof}

\paragraph{The logical open-cover step also requires an additional premise.}
There is a separate elementary reason that an open-cover or saturation
argument alone does not force a finite certificate pool. In $R^3$ let
$H=\{x_3=0\}$ and $V=R^3\setminus H$. Choose pairwise disjoint
open rectangles $I_n\subset H$, with rational endpoints, accumulating
only at the origin of $H$, and not containing it. For example take
\[
 I_n=(2^{-n},\tfrac54\,2^{-n})\times
                    (-2^{-n-3},2^{-n-3})\times\{0\},\qquad n\ge1.
\]
The sets
\[
 E_N=V\cup\bigcup_{n=1}^N I_n
\]
are increasing definable open sets. Their union $E$ is open and dense
and contains the fixed definable dense open set $V$. It is not
definable over the real field: its intersection with the positive
$x_1$-axis is an infinite union of separated intervals. Every point
of $E$ has a neighborhood inside one $E_N$, and all restrictions and
finite gluing of these neighborhoods remain valid. Nevertheless no
finite subfamily covers $E$.
Intersecting everything with the fixed compact cube $[-1,1]^3$ gives
the same example with relative open sets in a compact ambient space.
The set being covered remains a proper open subset of that cube, so
compactness of the ambient cube supplies no finite subcover of it.

The same formulas can be interpreted in every elementary extension;
every point of their external union still lies in one of its members.
This example is not asserted to be a chart locus of a definable map.
It isolates the precise logical failure: openness, a fixed definable
dense open subset, local witnesses, and finite gluing do not imply a
finite pool. Compactness would give a finite pool if the witness
condition covered every point of one fixed definable set in every
model, or if the exact union were already known to be definable in a
sufficiently saturated model, as in A5.019. That additional premise
is the unresolved recognition requirement.

\paragraph{Why the planar theorem cannot silently be made parameter-uniform.}
The complete recovered proof of A5.013 was reread for this diagnostic.
Its finite-format argument produces local PL models on the definable
punctured-edge set $D_e$. It then handles the finite complement of
each edge, and the finitely many target vertices, by observing that
any subset of a finite set is definable with parameters. That is valid
for the one fixed map in its statement. It supplies no formula for
choosing the good exceptional points as an external family of maps
varies. Even one such exceptional point for each parameter can carry
an undetermined parameter predicate. Consequently applying A5.013
separately to the positive-radius angular maps of A5.019 does not
establish the required family-uniform finite certificate pool.

\paragraph{Finite verification and scope.}
The exact symbolic checker
\path{agents/maximal_locus/code/check_hardt_stratum_fold.py}
verifies the square-circle joins and identified endpoint, fold and
homotopy endpoints, all affine slopes, and the complete three-point
inverse fibre in (HT.5)--(HT.6). Its actual saved output
\path{agents/maximal_locus/logs/check_hardt_stratum_fold.json}
records $31$ passes in $0.06948413099962636$ seconds. No numerical
sampling is used. The properness, nonapproximation, and local-bundle
claims are proved above and are not delegated to these finite checks.
This note is an unnumbered limitation of two proposed proof methods.
It supplies no nondefinable bundle locus and no movement of original A5.
